
\section{11. Conclusion}

\subsection{11.1 Summary}

We presented \textbf{Progressive Token Pruning} (PTP), a learnable token-selection 
module for efficient transformer attention. PTP scores each token with a 
small MLP, retains the top-$k$ fraction, and preserves [CLS] and causality.

\subsection{11.2 Results Recap}

\begin{table}[H]
\centering
\small
\begin{tabular}{llll}
\toprule
\textbf{Model} & \textbf{Acc $\Delta$} & \textbf{Speedup} & \textbf{Kept} \\
\midrule
DistilBERT & -1.2\% & 1.25$\times$ & 80\% \\
RoBERTa & -1.6\% & 1.20$\times$ & 80\% \\
GPT-2 & -1.4\% & 1.20$\times$ & 80\% \\
\bottomrule
\end{tabular}
\end{table}

\subsection{11.3 Key Takeaways}

1. \textbf{Token pruning is effective} at $\rho = 0.8$.
2. \textbf{[CLS] preservation is essential.}
3. \textbf{Learned scores reveal structure.}
4. \textbf{Causality preserved.}
5. \textbf{Drop-in wrapper.}

\subsection{11.4 Future Work}

\subsubsection{Short-Term}
\begin{itemize}
\item Causal score network
\item Per-head scoring
\item Learned keep ratio
\item Long-context evaluation
\end{itemize}

\subsubsection{Long-Term}
\begin{itemize}
\item Cross-layer score sharing
\item Theoretical bounds
\item Integration with MoE
\end{itemize}

\subsection{11.5 Closing Remarks}

PTP provides a practical, learnable, drop-in token-pruning module. We 
release code and checkpoints.

\subsection{11.6 Acknowledgments}

We thank {{ACKNOWLEDGMENTS}}.

\subsection{11.7 Author Contributions}

\begin{itemize}
\item \textbf{{{AUTHOR\_1}}} designed PTP, implemented the code, wrote the paper.
\item \textbf{{{AUTHOR\_2}}} provided feedback, edited.
\end{itemize}

\subsection{11.8 Conflict of Interest}

None.